Niniejsza praca dyplomowa zrealizowana została zgodnie z założonymi celami badawczymi. Przeprowadzono kompleksową analizę wydajnościową i funkcjonalną ośmiu wybranych nierelacyjnych baz danych typu klucz-wartość (Redis, Memcached, LevelDB, RocksDB, Couchbase, Tarantool, FoundationDB oraz BerkeleyDB) w środowisku natywnego systemu operacyjnego Linux (Ubuntu) oraz zwirtualizowanego WSL 2.

W ramach pracy zaprojektowano i zaimplementowano autorską aplikację benchmarkującą w języku Python. Narzędzie to umożliwiło przeprowadzenie zróżnicowanych testów wydajnościowych (operacje CRUD, obciążenia mieszane, operacje na strukturach JSON oraz kolejkach FIFO) przy precyzyjnym limicie zasobów sprzętowych (CPU, RAM) narzucanym przez mechanizm cgroups w kontenerach Docker. Zgromadzone dane pomiarowe poddano szczegółowej analizie statystycznej, wyznaczając średnie oraz odchylenia standardowe przepustowości, a także centyle opóźnień (p50, p95, p99), co pozwoliło na identyfikację stabilności badanych silników i opóźnień ogonowych.

Główne wnioski z przeprowadzonych badań można podsumować następująco:
\begin{itemize}
    \item \textbf{Wpływ architektury na wydajność:} Systemy operujące w całości w pamięci RAM (Memcached, Redis) wykazują najwyższą i najbardziej stabilną przepustowość w konfiguracjach klient-serwer. Biblioteki osadzone (RocksDB, LevelDB) dominują w środowiskach jednowątkowych, eliminując narzut sieciowy.
    \item \textbf{Narzut wirtualizacji WSL 2:} Badania wykazały, że nowoczesna wirtualizacja w WSL 2 charakteryzuje się minimalnym narzutem na komunikację sieciową loopback (do 8\% spadek przepustowości). Zidentyfikowano jednak znaczną degradację wydajności dyskowej (od 25\% do 40\%) dla baz transakcyjnych wymuszających synchroniczny zapis WAL, co jest związane z translacją wywołań \texttt{fsync} na wirtualnym dysku VHDX NTFS.
    \item \textbf{Weryfikacja wielokryterialna:} Końcowy ranking wskazał systemy Redis oraz Tarantool jako najbardziej optymalne, zrównoważone rozwiązania produkcyjne, które łączą bogatą natywną funkcjonalność z wysoką wydajnością. Wykazano również unikalne cechy systemu BerkeleyDB, będącej jedynym testowanym silnikiem wbudowanym umożliwiającym bezpieczną i efektywną pracę współbieżną z wielu procesów systemowych.
\end{itemize}

Jako kierunki dalszych prac rozwojowych w tym obszarze badawczym wskazuje się:
\begin{itemize}
    \item \textbf{Badania w środowisku rozproszonym (skalowanie poziome):} Przeprowadzenie analogicznych pomiarów dla systemów klastrowych z uwzględnieniem narzutu opóźnień sieci fizycznej (LAN/WAN) oraz protokołów spójności.
    \item \textbf{Optymalizacja kodu klienta:} Przepisanie aplikacji benchmarkującej na język o wyższej wydajności (np. Go lub Rust) w celu wyeliminowania wąskiego gardła jednowątkowego interpretera CPython przy generowaniu bardzo intensywnego ruchu współbieżnego.
    \item \textbf{Złożone profile obciążenia:} Zaimplementowanie generatora obciążeń o rozkładzie niesymetrycznym (np. rozkład Zipfa) w celu dokładniejszego odwzorowania zachowania baz danych w rzeczywistych systemach produkcyjnych.
\end{itemize}